%%% Please fill in basic information on your thesis, which will be automatically
%%% inserted at the right places. You need to replace \xxx{...} by real data.

% Type of your thesis:
%	"bc" for Bachelor's
%	"mgr" for Master's
%	"phd" for PhD
%	"rig" for rigorosum
\def\ThesisType{mgr}

% Language of your study programme:
%	"cs" for Czech
%	"en" for English
\def\StudyLanguage{en}

% Thesis title in English (exactly as in the official assignment)
% (Note: \xxx is a "ToDo label" which makes the unfilled visible. Remove it.)
\def\ThesisTitle{Evolutionary Algorithms for Project Scheduling}

% Author of the thesis (you)
\def\ThesisAuthor{Sabína Ságová}

% Year when the thesis is submitted
\def\YearSubmitted{2026}

% Name of the department or institute, where the work was officially assigned
% (according to the Organizational Structure of MFF UK in English,
% see https://www.mff.cuni.cz/en/faculty/organizational-structure,
% or a full name of a department outside MFF)
\def\Department{Department of Theoretical Computer Science and Mathematical Logic}

% Is it a department (katedra), or an institute (ústav)?
\def\DeptType{Department}

% Thesis supervisor: name, surname and titles
\def\Supervisor{doc. Mgr. Martin Pilát, Ph.D.}

% Supervisor's department (again according to Organizational structure of MFF)
\def\SupervisorsDepartment{Department of Theoretical Computer Science and Mathematical Logic}

% Study programme (does not apply to rigorosum theses)
\def\StudyProgramme{Computer Science - Artificial Intelligence}

% An optional dedication: you can thank whomever you wish (your supervisor,
% consultant, who provided you with tea and pizza, etc.)
\def\Dedication{%
\vspace*{1cm}

\noindent
First and foremost, I would like to thank my supervisor, \Supervisor, for
agreeing to supervise this project, and for the supportive environment and
steady guidance he provided throughout its development.

I am also sincerely grateful to Yuan Tian, PhD candidate at Victoria
University of Wellington, for generously providing the source code of the
baseline algorithm used throughout this thesis and for patiently taking the
time to explain it to me.

Computational resources for the experiments in this thesis were provided
by the e-INFRA CZ project (ID:90254), supported by the Ministry of
Education, Youth and Sports of the Czech Republic; I thank MetaCentrum for
access to its computing grid, without which the full-scale experimental
evaluation would not have been possible.

On a more personal note, I want to thank my father, and remember my
grandfather with love, for always believing in me and for making me
independent and confident in myself.

Finally, thanks to my friends Andrea and Kateryna for their support along
the way, and for sharing these university years with me.
}

\def\AIAssistance{\subsection*{Disclosure statement} In accordance with \textbf{Dean’s Directive No. 9/2024} (amending the Code of Ethics of the Faculty of Mathematics and Physics, Charles University), I disclose that \textbf{Grammarly, GitHub Copilot, and Gemini 3 Flash} were used as supportive tools for language polishing, grammatical refinement, and code assistance. 
These tools were not used as sources of information, nor did they replace original critical analysis or literature review. As the author, I bear \textbf{sole responsibility} for the accuracy and ethical integrity of this work, and I affirm that these AI tools are not considered co-authors.
}
% Abstract (recommended length around 80-200 words; this is not a copy of your thesis assignment!)
\def\Abstract{%
This thesis extends a genetic-programming hyper-heuristic (GPHH) baseline
for the strongly NP-hard multi-mode resource-constrained project
scheduling problem. A staged methodology first screens a wide field of
candidate modifications cheaply, then evaluates the survivors at full
statistical power on MMLIB50 and MMLIB100. \changed{Of the screened variants,
including a diagnosis-driven graft-based variation scheme, only one,
epsilon-lexicase selection,} together with two independently motivated
extensions, critical-path local search and non-renewable (NR) resource
terminals, is carried forward. NR terminals, which expose previously
invisible non-renewable resource state to the evolving rule, deliver a
large, statistically significant improvement in \changedB{schedule quality
that replicates robustly across benchmark scale and on the complete
held-out test set; the accompanying feasibility gain proves largely
confined to instance classes seen during training.} Lexicase
selection reliably improves training fitness but this advantage does not
transfer reliably to held-out test fitness, and the gap widens rather than
narrows at larger scale. \changed{Critical-path local search improves schedule
quality on the smaller benchmark, but only under one of the two schedule
generation schemes, and the improvement does not survive the move to the
larger one.}
}

% 3 to 5 keywords (recommended) separated by \sep
% Keywords are useful for indexing and searching for the theses by topic.
\def\ThesisKeywords{%
multi-mode project scheduling\sep
genetic programming\sep
hyper-heuristics\sep
evolutionary algorithms}

% If any of your metadata strings contains TeX macros, you need to provide
% a plain-text version for use in XMP metadata embedded in the output PDF file.
% If you are not sure, check the generated thesis.xmpdata file.
\def\ThesisAuthorXMP{\ThesisAuthor}
\def\ThesisTitleXMP{\ThesisTitle}
\def\ThesisKeywordsXMP{\ThesisKeywords}
% Plain-text abstract for XMP metadata (the \Abstract above contains the
% \changed revision-highlighting macro, which must not leak into metadata).
\def\AbstractXMP{This thesis extends a genetic-programming hyper-heuristic
(GPHH) baseline for the strongly NP-hard multi-mode resource-constrained
project scheduling problem. A staged methodology first screens a wide field
of candidate modifications cheaply, then evaluates the survivors at full
statistical power on MMLIB50 and MMLIB100. Of the screened variants,
including a diagnosis-driven graft-based variation scheme, only one,
epsilon-lexicase selection, together with two independently motivated
extensions, critical-path local search and non-renewable (NR) resource
terminals, is carried forward. NR terminals, which expose previously
invisible non-renewable resource state to the evolving rule, deliver a
large, statistically significant improvement in schedule quality that
replicates robustly across benchmark scale and on the complete held-out
test set; the accompanying feasibility gain proves largely confined to
instance classes seen during training. Lexicase
selection reliably improves training fitness but this advantage does not
transfer reliably to held-out test fitness, and the gap widens rather than
narrows at larger scale. Critical-path local search improves schedule
quality on the smaller benchmark, but only under one of the two schedule
generation schemes, and the improvement does not survive the move to the
larger one.}

% If your abstracts are long and do not fit in the infopage, you can make the
% fonts a bit smaller by this setting. (Also, you should try to compress your abstract more.)
\def\InfoPageFont{}
%\def\InfoPageFont{\small}  % uncomment to decrease font size

% If you are studing in a Czech programme, you also need to provide metadata in Czech:
% (in English programmes, this is not used anywhere)

\def\ThesisTitleCS{Evoluční algoritmy pro rozvrhování projektů}
\def\DepartmentCS{Katedra teoretické informatiky a matematické logiky}
\def\DeptTypeCS{Katedra}
\def\SupervisorsDepartmentCS{Katedra teoretické informatiky a matematické logiky}
\def\StudyProgrammeCS{Informatika -- umělá inteligence}

\def\ThesisKeywordsCS{%
více-režimové rozvrhování projektů\sep
genetické programování\sep
hyperheuristiky\sep
evoluční algoritmy
}

\def\AbstractCS{%
Tato práce rozšiřuje hyperheuristiku založenou na genetickém programování
(GPHH) pro silně NP-těžký problém rozvrhování projektů s omezenými zdroji
a více režimy provádění aktivit (MRCPSP). Vícestupňová metodika nejprve
levně prověří širokou množinu kandidátních modifikací a přeživší kandidáty
poté vyhodnotí s plnou statistickou silou na benchmarcích MMLIB50 a
MMLIB100. \changed{Z prověřených variant, včetně diagnosticky řízeného
schématu variace založeného na roubování podstromů, je dále rozvinuta}
pouze epsilon-lexicase selekce, spolu se dvěma nezávisle motivovanými
rozšířeními: lokálním prohledáváním řízeným kritickou cestou a terminály
pro neobnovitelné zdroje (NR). NR terminály, které evolvovanému pravidlu
zpřístupňují dosud neviditelný stav neobnovitelných zdrojů, přinášejí
velké a statisticky významné zlepšení \changedB{kvality rozvrhů, které se
robustně replikuje napříč velikostmi benchmarku i na úplné testovací sadě;
doprovodné zlepšení přípustnosti se však ukazuje být převážně omezeno na
třídy instancí viděné při trénování.}
Lexicase selekce spolehlivě zlepšuje trénovací fitness, tato výhoda se
však nepřenáší na testovací data a s rostoucí velikostí instancí se rozdíl
dále prohlubuje. \changed{Lokální prohledávání řízené kritickou cestou zlepšuje
kvalitu rozvrhů na menším benchmarku, ale pouze pod jedním ze dvou schémat
generování rozvrhu, a toto zlepšení se na větším benchmarku nereplikuje.}
}
